% Katalog A: Struktur und Navigation (Palette Slot 1).
\ZSFNewColumn
\StartChapter[kat:struktur]{A · Struktur}

\begin{runintext}
  Struktur-Makros haben keinen Inhalt — ihr Katalogeintrag ist der Balken
  selbst. Die Kapitelzeile über diesem Absatz ist A-02, jeder Abschnittsbalken
  unten sein eigener Eintrag.
\end{runintext}

\begin{defbox}[A-01 · ZSFTitleHeader][weight=quiet]
  \ZSFindex{ZSFTitleHeader}
  Steht genau einmal ganz oben im Dokument — in dieser Ausgabe über der
  Lesehilfe. Nicht wiederholbar, deshalb hier nur benannt.
\end{defbox}
\Zweck{Dokumentkopf mit Titel und Autor.}

\begin{defbox}[A-02 · StartChapter][weight=quiet]
  \ZSFindex{StartChapter}
  Der breite Balken am Kopf dieser Sektion. Nummeriert automatisch und
  vergibt die Kapitelfarbe an alles, was folgt.
\end{defbox}
\Zweck{Hauptkapitel mit Auto-Nummer und eigener Palette-Farbe.}

\begin{defbox}[A-03 · StartFrontChapter][weight=quiet]
  \ZSFindex{StartFrontChapter}
  Unnummeriert und in ETH-Blau — der Balken über der Lesehilfe. Das
  Kurzlabel im optionalen Argument dient dem Register als Wegweiser.
\end{defbox}
\Zweck{Front- und Back-Matter ohne Kapitelnummer.}

\SubsectionBar[kat:bar]{A-04 · SubsectionBar}
\ZSFindex{SubsectionBar}

\begin{defbox}[][weight=quiet]
  \Muster
\end{defbox}
\Zweck{Nummerierter Abschnitt; liefert Label für Verweise und Index-Locator.}

\SubsectionBar*{A-05 · SubsectionBar*}

\begin{defbox}[][weight=quiet]
  \Muster
\end{defbox}
\Zweck{Gliederung ohne neue Nummer — rein visuell, kein Verweisziel.}

\SubsubsectionBar[kat:subsub]{A-06 · SubsubsectionBar}
\ZSFindex{SubsubsectionBar}

\begin{defbox}[][weight=quiet]
  \Muster
\end{defbox}
\Zweck{Dritte Lookup-Ebene, sichtbar kleiner und heller.}

\SubsubsectionBar*{A-07 · SubsubsectionBar*}

\begin{defbox}[][weight=quiet]
  \Muster
\end{defbox}
\Zweck{Dritte Ebene ohne Nummer.}

\SubsectionBar*{A-08 · ZSFNewColumn}

\begin{defbox}[][weight=quiet]
  \ZSFindex{ZSFNewColumn}
  Unsichtbar: erzwingt den Spaltenwechsel vor einem Strukturmakro. Steht in
  diesem Katalog vor jedem Sektionskopf, damit jede Sektion oben beginnt.
\end{defbox}
\Zweck{Bewusster Spaltenumbruch statt rohem columnbreak.}
